\subsection{Static Analysis Before Deployment}
\label{sec:predeploy}

Runtime mediation decides calls as they happen; a static scan decides
what an agent \emph{could} do before it runs, and the two share a policy
vocabulary. The pre-deployment scanner parses an agent's repository with
an AST and taint-tracking scanner (we integrate agent-audit rather than
maintain a second regex scanner), enumerates every model call, tool
definition and MCP server configuration, and applies agent-audit's rules (the version the gateway recommends is pinned in \texttt{predeploy-scan.ts}), mapped to the OWASP Agentic Top~10, so a finding carries the same category an
auditor will ask about. Scans are kept as history, and a scan-to-scan
diff partitions findings into \emph{added}, \emph{fixed} and
\emph{persisting}, which is the form in which a reviewer wants to see a
pull request: not ``this agent has twelve findings'' but ``this change
introduced one''. A code shield applies a small set of high-precision
rules to code an agent is about to commit, execute or hand back, in the
spirit of LlamaFirewall's CodeShield but kept deliberately narrow; and
the runtime's L1 patterns include a supply-chain family (source-map leak,
unsafe publish, secret in a build artefact) that fires on the publish and
build tool calls an agent makes, so the static and dynamic views meet
at the same boundary.
